% !TeX root = qc_report.tex
% !TeX program = pdflatex
\documentclass[11pt,a4paper]{article}

\pdfminorversion=7

% =====================
% Packages
% =====================
\usepackage[top=0.9in,left=0.9in,right=0.9in,bottom=1.2in]{geometry}
\usepackage{array}
\usepackage{ragged2e}
\usepackage{setspace}
\usepackage{hyperref}
\usepackage{enumitem}
\usepackage{graphicx}
\usepackage{fancyhdr}
\usepackage{tabularx}
\usepackage{booktabs}
\usepackage{multirow}
\usepackage{float}
\usepackage{xcolor}
\usepackage{parskip}
\usepackage{amsmath}
\usepackage{amssymb}
\usepackage{colortbl}
\usepackage{longtable}
\usepackage{textcomp}
\usepackage{caption}
\usepackage{bookmark}
\usepackage[normalem]{ulem}

\hypersetup{
colorlinks=true,
linkcolor=blue,
urlcolor=blue,
citecolor=blue
}

\AtBeginDocument{
\let\origref\ref
\renewcommand{\ref}[1]{\uline{\origref{#1}}}
\let\orighyperref\hyperref
\renewcommand{\hyperref}[2][]{\orighyperref[#1]{\uline{#2}}}
}

\definecolor{tableheader}{RGB}{240,240,240}
\definecolor{tablerow}{RGB}{250,250,250}

\frenchspacing

% =====================
% Font
% =====================
\usepackage[T1]{fontenc}
\usepackage[scaled=0.95]{helvet}
\renewcommand{\familydefault}{\sfdefault}

\setstretch{1.0}

% =====================
% HEADER/FOOTER
% =====================
\pagestyle{fancy}
\fancyhf{}
\lhead{\textbf{Veridion Systems Validation Labs}}
\rhead{Quality Control (QC) Audit Report}
\cfoot{\thepage}
\renewcommand{\headrulewidth}{0.4pt}
\setlength{\headheight}{20pt}

\begin{document}

% =====================
% COVER PAGE
% =====================
\thispagestyle{empty}
\fancyhf{}
\renewcommand{\headrulewidth}{0pt}

\newgeometry{top=5.2cm,left=0.9in,right=0.9in,bottom=1.2in}

\begin{center}

{\Large\fbox{\parbox{0.72\textwidth}{\centering
\vspace{0.3cm}
\textbf{BACKGROUND PLACEHOLDER}\\
Insert QC cover background here later
\vspace{0.3cm}}}}

\vspace{1cm}

\LARGE\textbf{Quality Control (QC) Audit Report}

\large Software Requirements Specification (SRS)

\large Smart Contract Enabled Pharmaceutical Supply Chain System

\vspace{0.8cm}

\textbf{\Large Veridion Systems Validation Labs}

\today

\end{center}

\begin{table}[H]
\centering
\small
\renewcommand{\arraystretch}{2.0}
\begin{tabular}{|m{0.25\textwidth}|m{0.2\textwidth}|m{0.25\textwidth}|m{0.2\textwidth}|}
\hline
\rowcolor{tableheader}
Brigman Team Member & Department & Role & Employee ID \\
\hline
Muhammad Ridwan Putra Jasni & QC Team & QC Audit Lead & QC-001 \\
\hline
Lin Yuhao & QC Team & QC Specialist & QC-002 \\
\hline
Muhammad Mikhail Bin Mazlan & QC Team & QC Specialist & QC-003 \\
\hline
Kannan s/o Rajamohan & QC Team & QC Specialist & QC-004 \\
\hline
\end{tabular}
\caption{QC Team -- Veridion Systems Validation Labs}
\label{tab:qc-team}
\end{table}

\restoregeometry
\newpage

% =====================
% TABLE OF CONTENTS
% =====================
\renewcommand{\contentsname}{Table of Contents}
\thispagestyle{empty}
\hypersetup{linkcolor=black}
\tableofcontents
\newpage
\hypersetup{linkcolor=blue}

\noindent\fbox{\begin{minipage}{\dimexpr\textwidth-2\fboxsep-2\fboxrule\relax}
\small\textbf{Note on Hyperlinks:} Throughout this document, all hyperlinks are displayed in \textcolor{blue}{\underline{blue and underlined}}. These links are fully clickable in the PDF viewer and will navigate to the relevant section, finding, or appendix.
\end{minipage}}

\bigskip

% =====================
% EXECUTIVE SUMMARY
% =====================
\section{Executive Summary}

This report presents the results of a Quality Control (QC) validation of the Software Requirements Specification (SRS) for the \textbf{Smart Contract Enabled Pharmaceutical Supply Chain System}. The objective of this QC audit is to identify defects that remain in the SRS after the completion of Eliciting, Analysing, and Specifying.

The audit focused on three principal categories of SRS defects expected at the Validating stage:
\begin{itemize}[leftmargin=*]
\item \textbf{Authorship defects} -- where intended meanings are not expressed clearly enough.
\item \textbf{Omission defects} -- where required scenarios, conditions, or supporting content are missing.
\item \textbf{Contradiction defects} -- where different parts of the SRS imply inconsistent behaviour.
\end{itemize}

The inspection found that the SRS is generally well-structured and covers most core functional areas. However, a number of issues were detected including contradictory requirements (MFA scope, blockchain confirmation time), missing exception-handling scenarios (blockchain failure, invalid QR codes, password recovery), vague statements, and typographical errors. A glossary is absent, and some cross-references are inconsistent.

These findings indicate that the SRS is \textbf{conditionally acceptable}, subject to the customer's review of the identified defects. No recommendations for correction are provided; the customer will determine any necessary actions.

% =====================
% 1 INTRODUCTION
% =====================
\section{Introduction}

\subsection{Purpose}
This Quality Control (QC) Audit Report presents the findings of an independent inspection of the Software Requirements Specification (SRS) prepared for the \textbf{Smart Contract Enabled Pharmaceutical Supply Chain System}. The purpose of this QC activity is to identify defects remaining in the SRS work product before the project progresses into downstream stages such as system design and implementation.

In accordance with the validating stage of Requirements Engineering, QC focuses on defect detection in the \textbf{requirements work product itself}. The QC team records findings but does not propose fixes; the customer determines any corrective actions.

\subsection{Document Conventions}
This document is an audit report rather than a replacement SRS. It records detected defects and explains their implications.

Audit findings use the format \texttt{QC-FND-XXX}.

Error detection instances use the format \texttt{QC-ERR-XXX}.

\subsection{Intended Audience and Reading Suggestions}
This report is intended for the Medicare Plus requirements engineering team, project management personnel, and relevant stakeholders involved in SRS review and approval.

Readers should begin with the Executive Summary and Sections 1--3, then proceed to the QC Checklist, Error Detection, and detailed Findings.

\subsection{Audit Scope}
The audit covers the SRS as a QC exercise. The scope includes:
\begin{itemize}[leftmargin=*]
\item detection of authorship defects (ambiguity, vagueness, grammatical errors)
\item detection of latent contradictions
\item detection of omission of topics (missing scenarios, exception paths)
\item checking codification and traceability coherence
\item checking structural, editorial, and presentation consistency
\end{itemize}

% =====================
% 2 AUDIT BASIS
% =====================
\section{Audit Basis}

\subsection{Basis of Evaluation}
This QC audit was conducted on the understanding that the SRS should already be in reasonably strong condition after completion of Eliciting, Analysing, and Specifying. The purpose of QC at this stage is therefore not to rewrite the SRS, but to detect remaining defects still likely to be present in the work product.

\subsection{Nature of QC}
Quality Control focuses on inspecting the SRS work product itself. The QC team identifies defects but does not directly amend the SRS. Instead, the audit team produces a report that identifies defects for consideration by the customer.

\subsection{Kinds of Errors Pursued}
The QC audit focuses on detecting the following major defect categories expected in the SRS at the validating stage:
\begin{itemize}[leftmargin=*]
\item \textbf{Indiscriminate side effects of authorship} -- where intended meaning is clear to the writers but not expressed with sufficient precision.
\item \textbf{Latent contradictions} -- where requirements, models, or sections conflict with one another.
\item \textbf{Omission of topics} -- where expected requirements content, edge conditions, or supporting information is absent.
\item \textbf{Traceability coherence issues} -- where codification, source mapping, or cross-references are incomplete or unclear.
\item \textbf{Editorial and structural inconsistencies} -- where formatting, section naming, or presentation patterns weaken readability or professionalism.
\end{itemize}

\subsection{Severity Classification}
Findings are categorised according to severity:
\begin{itemize}[leftmargin=*]
\item \textbf{High} -- defect that may significantly affect interpretation, implementation, or expected system behaviour.
\item \textbf{Medium} -- defect that may introduce ambiguity or reduce clarity without fundamentally changing requirement intent.
\item \textbf{Low} -- minor editorial or structural issue that does not materially affect requirement meaning.
\end{itemize}

\subsection{Documents Reviewed}
The audit considered the following project materials as supporting references:
\begin{itemize}[leftmargin=*]
\item Software Requirements Specification (SRS) for the Smart Contract Enabled Pharmaceutical Supply Chain System, version as of 6 March 2026.
\item Project Specification (assumed background).
\item Relevant requirements models (context diagram, sequence diagram, architecture diagram).
\item Codification and iRTM sections.
\item Quality model and change mitigation sections.
\item Appendices and supporting traceability artefacts.
\end{itemize}

% =====================
% 3 QC PLAN
% =====================
\section{QC Plan}

\subsection{Target Workproduct}
The target workproduct for this QC audit is the released \textbf{Software Requirements Specification (SRS)} for the Smart Contract Enabled Pharmaceutical Supply Chain System project.

\subsection{Context and Goals}
The SRS is the main requirements work product produced after Analysing and Specifying. Since downstream design and implementation rely on the SRS as the controlling requirements baseline, the QC audit is conducted to provide confidence that the specification is sufficiently clear, complete, and internally consistent.

The goals of this QC audit are:
\begin{itemize}[leftmargin=*]
\item to determine whether the SRS clearly and consistently represents the intended requirements.
\item to identify material defects before downstream design and implementation begin.
\item to document defects for customer consideration.
\end{itemize}

\subsection{Access and Audit Conditions}
The QC team reviewed the released SRS and supporting project artefacts made available for audit. Access was limited to materials needed to inspect the work product and its supporting traceability and modelling references.

\subsection{Inspection Techniques Used}
The following investigative techniques, drawn from industry best practices and the lecture material, were applied during the QC audit:
\begin{enumerate}[leftmargin=*]
\item \textbf{Close document inspection} -- detailed reading of requirement clauses and supporting sections.
\item \textbf{Perspective-based reading} -- reviewing the SRS from likely reader viewpoints such as developer, tester, project manager, and compliance officer.
\item \textbf{Cross-section contradiction checking} -- comparing different SRS sections and models to detect hidden conflicts (e.g., performance targets, security scopes).
\item \textbf{Codification and traceability verification} -- checking requirement IDs, source links, iRTM mappings, and cross-references.
\item \textbf{Model-to-text consistency inspection} -- ensuring models (diagrams) and textual requirements express compatible meanings.
\item \textbf{Omission mapping against expected SRS topics} -- checking whether important topics, edge cases, and exception paths are sufficiently covered (e.g., failure handling, recovery processes).
\item \textbf{Scenario-based reading} -- testing the SRS against realistic supply-chain scenarios (order placement, product verification, blockchain outage).
\item \textbf{Assumption reading} -- identifying assumptions that contradict documented requirements or have been rendered redundant.
\item \textbf{Double-checking} -- verifying dates, facts, figures, and proper names for consistency.
\item \textbf{Spellcheck and grammar review} -- using automated tools plus manual reading to catch typographical errors.
\end{enumerate}

\subsection{Work Breakdown}
The QC audit team consisted of four members, with one member also serving as the QA auditor (not reflected in this QC report). The work was allocated as follows:

\begin{table}[H]
\centering
\small
\renewcommand{\arraystretch}{1.4}
\begin{tabularx}{\textwidth}{|p{0.22\textwidth}|p{0.28\textwidth}|X|}
\hline
\rowcolor{tableheader}
\textbf{Role} & \textbf{Assigned Member} & \textbf{Primary Responsibility} \\
\hline
QC Audit Lead & Muhammad Ridwan Putra Jasni & Coordinate audit activities, consolidate findings, ensure consistency of final report, apply cross-section contradiction checking and assumption reading. \\
\hline
Requirements Reviewer & Lin Yuhao & Inspect operational, technology, and quality-related sections for ambiguity, contradiction, and omission; perform perspective-based reading (developer, tester) and scenario-based reading. \\
\hline
Traceability Reviewer & Muhammad Mikhail Bin Mazlan & Review codification, iRTM references, source mappings, and section-to-section coherence; double-check all identifiers and cross-references. \\
\hline
Editorial and Structure Reviewer & Kannan s/o Rajamohan & Review document structure, section naming, formatting consistency, and overall presentation quality; conduct spellcheck and grammar review, verify figures and tables. \\
\hline
\end{tabularx}
\caption{QC audit team work allocation}
\label{tab:qc-wbs}
\end{table}

\subsection{Indicative Timeline}
The audit was conducted over a single day, with the following estimated timeline:

\begin{table}[H]
\centering
\small
\renewcommand{\arraystretch}{1.4}
\begin{tabularx}{\textwidth}{|p{0.18\textwidth}|p{0.42\textwidth}|X|}
\hline
\rowcolor{tableheader}
\textbf{Stage} & \textbf{Activity} & \textbf{Indicative Duration} \\
\hline
Stage 1 & Individual inspection of assigned SRS sections using selected techniques (close reading, perspective-based, scenario-based) & 3 hours \\
\hline
Stage 2 & Consolidation meeting and de-duplication of identified findings & 1 hour \\
\hline
Stage 3 & Severity review and documentation of findings & 1 hour \\
\hline
Stage 4 & Final report formatting, checking, and team sign-off & 1 hour \\
\hline
\end{tabularx}
\caption{Indicative QC audit timeline}
\label{tab:qc-timeline}
\end{table}

\subsection{Reporting}
The QC team records findings in this report. The audit team does not directly amend the audited SRS.

% =====================
% 4 QC CHECKLIST
% =====================
\section{QC Checklist}

The following checklist was used by the QC audit team to inspect the SRS work product. The checklist items represent critical aspects of a high-quality SRS, derived from industry standards and lecture guidelines.

\begingroup
\setlength{\LTpre}{0pt}
\setlength{\LTpost}{0pt}
\setlength{\tabcolsep}{4pt}
\small
\renewcommand{\arraystretch}{1.35}
\begin{longtable}{|>{\RaggedRight\arraybackslash}p{0.08\textwidth}|>{\RaggedRight\arraybackslash}p{0.14\textwidth}|>{\RaggedRight\arraybackslash}p{0.56\textwidth}|>{\RaggedRight\arraybackslash}p{0.16\textwidth}|}
\caption{QC inspection checklist}\label{tab:qc-checklist}\\
\hline
\rowcolor{tableheader}
\textbf{Item} & \textbf{Focus} & \textbf{QC Checklist Item} & \textbf{\shortstack[l]{Inspection\\Result}} \\
\hline
\endfirsthead
\hline
\rowcolor{tableheader}
\textbf{Item} & \textbf{Focus} & \textbf{QC Checklist Item} & \textbf{\shortstack[l]{Inspection\\Result}} \\
\hline
\endhead
\hline \multicolumn{4}{r}{{Continued on next page}} \\
\endfoot
\hline
\endlastfoot

\rowcolor{tablerow}\multicolumn{4}{|l|}{\textbf{1. Introduction Section}} \\
\hline
QC1 & Completeness & Does the Introduction explicitly define the system purpose, target user base, and operational boundaries (in-scope vs. out-of-scope)? & Satisfied \\
\hline
QC2 & Completeness & Does the Introduction state the specific business problem and identify measurable business goals or KPIs? & Partially satisfied \\
\hline
QC3 & Consistency & Is the system name, project title, and version number identical across title page, headers/footers, and body? & Satisfied \\
\hline

\rowcolor{tablerow}\multicolumn{4}{|l|}{\textbf{2. Glossary / Definitions Section}} \\
\hline
QC4 & Completeness & Are all domain-specific terms, acronyms, and technical abbreviations defined in a Glossary? & Not satisfied (missing) \\
\hline
QC5 & Relevance & Does every term defined in the Glossary appear at least once in the body? & Not applicable \\
\hline
QC6 & Consistency & Is each system role or entity referred to consistently using the same term throughout? & Satisfied \\
\hline

\rowcolor{tablerow}\multicolumn{4}{|l|}{\textbf{3. Overall Description Section}} \\
\hline
QC7 & Completeness & Are all user classes identified with their technical expertise, access privileges, and responsibilities? & Satisfied \\
\hline
QC8 & Completeness & Does the SRS specify the intended operational environment (OS, browsers, hardware, network)? & Satisfied \\
\hline
QC9 & Consistency & Are all listed dependencies and constraints adhered to and not contradicted by later requirements? & Partially satisfied \\
\hline

\rowcolor{tablerow}\multicolumn{4}{|l|}{\textbf{4. Functional Requirements Section}} \\
\hline
QC10 & Correctness & Does each requirement describe a single behaviour using precise, testable language (avoiding vague terms)? & Partially satisfied \\
\hline
QC11 & Testability & Does each functional requirement define trigger inputs, expected behaviour, and output/error state? & Partially satisfied \\
\hline
QC12 & Consistency & Is every requirement assigned a unique identifier following a consistent naming convention? & Satisfied \\
\hline

\rowcolor{tablerow}\multicolumn{4}{|l|}{\textbf{5. Interface Requirements Section}} \\
\hline
QC13 & Completeness & Does the Interface section describe purpose, data exchanged, and communication method for each external interface? & Partially satisfied \\
\hline
QC14 & Correctness & Is the trigger condition, expected data payload, and error handling defined for each interface? & Partially satisfied \\
\hline
QC15 & Consistency & Does every external interface listed correspond to at least one functional requirement? & Satisfied \\
\hline

\rowcolor{tablerow}\multicolumn{4}{|l|}{\textbf{6. Non-Functional Requirements Section}} \\
\hline
QC16 & Testability & Are all NFRs quantified with measurable limits (e.g., response time $\leq$ 2 sec for 95\% of queries)? & Partially satisfied \\
\hline
QC17 & Completeness & Does the SRS address all relevant NFR categories (performance, security, reliability, usability, maintainability)? & Satisfied \\
\hline
QC18 & Consistency & Do the NFRs logically align with system constraints (e.g., storage capacity consistent with data volume)? & Satisfied \\
\hline

\rowcolor{tablerow}\multicolumn{4}{|l|}{\textbf{7. Addenda / Supporting Artifacts}} \\
\hline
QC19 & Traceability & If an iRTM is included, does it map requirements to business goals and validation artifacts? & Satisfied \\
\hline
QC20 & Correctness & Does every diagram include a descriptive title, legend, and standard notation accurately representing the system? & Partially satisfied \\
\hline
QC21 & Consistency & Do all cross-references, figure numbers, and table citations point to correct existing sections/artifacts? & Partially satisfied \\
\hline

\end{longtable}
\endgroup

% =====================
% 5 ERROR DETECTION
% =====================
\section{Error Detection}

The QC audit identified instances of three principal QC error types: authorship, omission, and contradiction. Each detected instance is recorded with its location and characteristic.

\vspace{1\baselineskip}
\setcounter{table}{4}
\begingroup
\setlength{\LTpre}{0pt}
\setlength{\LTpost}{0pt}
\setlength{\tabcolsep}{4pt}
\renewcommand{\arraystretch}{1.4}
\begin{longtable}{|>{\RaggedRight\arraybackslash}p{0.15\textwidth}|>{\RaggedRight\arraybackslash}p{0.18\textwidth}|>{\RaggedRight\arraybackslash}p{0.25\textwidth}|>{\RaggedRight\arraybackslash}p{0.42\textwidth}|}
\hline
\rowcolor{tableheader}
\textbf{Error ID} & \textbf{Error Type} & \textbf{Location} & \textbf{Characteristic of Error} \\
\hline
\endfirsthead
\hline
\rowcolor{tableheader}
\textbf{Error ID} & \textbf{Error Type} & \textbf{Location} & \textbf{Characteristic of Error} \\
\hline
\endhead
\hline \multicolumn{4}{r}{{Continued on next page}} \\
\endfoot
\hline
\endlastfoot

QC-ERR-001 & Contradiction & Section 1.4.1 (In-Scope Deliverables, Authentication \& Security) vs Section 3.2 FR-2.6 & Section 1.4.1 requires MFA for "all user accounts", while FR-2.6 requires MFA only for "all administrative accounts". These statements directly conflict on the scope of MFA enforcement. \\
\hline
QC-ERR-002 & Authorship & Section 5.1 NFR-4 and NFR-7 & Duplicate word "shall" appears in both requirements: "The system shall shall complete 95\% of API responses within 2 seconds" and similar for database queries. This is a typographical error. \\
\hline
QC-ERR-003 & Omission & Section 3.2 (User \& Role Management) & FR-2.5 specifies account lockout after failed attempts, but no requirement exists for password reset or account recovery flow. This is a missing functional capability. \\
\hline
QC-ERR-004 & Authorship & Section 3.3 FR-3.9 & The term "periodic" is used without definition: "submit periodic GPS updates". The frequency or condition for updates is not specified, making the requirement ambiguous. \\
\hline
QC-ERR-005 & Omission & Section 3.1 (Blockchain Infrastructure) & FR-1.1 through FR-1.4 describe recording transactions on the blockchain, but there is no requirement addressing failure scenarios (e.g., blockchain network unavailable, transaction timeout, consensus failure). \\
\hline
QC-ERR-006 & Authorship & Section 1.4.3 System Boundaries & The section number 1.4.3 appears twice (pages 9 and 10) with different content, indicating a numbering or structural error. \\
\hline
QC-ERR-007 & Omission & Section 3.4 (Product Verification) & FR-4.1 allows customers to retrieve product lifecycle trace, but no requirement specifies system behavior when a scanned QR code is invalid, tampered, or not found in the blockchain. \\
\hline
QC-ERR-008 & Contradiction & Section 1.4.1 (Performance Requirements) vs Section 5.1 NFR-2 & Section 1.4.1 specifies blockchain transaction confirmation time as "<30 seconds for standard transactions". NFR-2 requires confirmation "within 10 seconds for standard operations". These targets are inconsistent. \\
\hline
QC-ERR-009 & Authorship & Section 5.5.3 QF-3 Tamper Resistance, Measurable Criteria & The criterion "100\% enforcement of role-based access control" is not measurable as written; it lacks a testable definition (e.g., how enforcement is verified). \\
\hline
QC-ERR-010 & Omission & Section 6.5 (Migration \& Transition) & While migration tools are mentioned, there are no functional requirements detailing the data migration process, validation of migrated data, or handling of migration failures. \\
\hline
QC-ERR-011 & Contradiction & Section 5.4 NFR-19 vs Section 6.2.5 CM-5 & NFR-19 specifies backup RPO $\leq$ 24 hours and RTO $\leq$ 4 hours. CM-5 states RTO $\leq$ 2 hours and RPO $\leq$ 30 minutes. These recovery objectives conflict. \\
\hline
QC-ERR-012 & Omission & Entire document & The SRS lacks a glossary or definitions section, leaving many acronyms (HSA, MOH, PDPC, etc.) undefined within the document itself. \\
\hline
\end{longtable}
\endgroup
\begin{center}
\addtocounter{table}{-1}
\captionof{table}{Detected QC error instances}
\label{tab:qc-error-detection}
\end{center}

% =====================
% 6 QC FINDINGS
% =====================
\section{QC Findings}

\subsection{Finding Summary}

\begin{table}[H]
\centering
\small
\renewcommand{\arraystretch}{1.4}
\begin{tabularx}{\textwidth}{|p{0.15\textwidth}|p{0.22\textwidth}|p{0.15\textwidth}|X|}
\hline
\rowcolor{tableheader}
\textbf{Finding ID} & \textbf{Type} & \textbf{Severity} & \textbf{Short Description} \\
\hline
QC-FND-001 & Contradiction & High & Conflicting statements about MFA requirements (all users vs. only administrators). \\
\hline
QC-FND-002 & Contradiction & High & Inconsistent blockchain confirmation time targets (30 seconds vs. 10 seconds). \\
\hline
QC-FND-003 & Contradiction & High & Conflicting RPO/RTO values between NFR-19 and CM-5. \\
\hline
QC-FND-004 & Omission & High & Missing exception handling for blockchain failures, invalid QR codes, password recovery, and migration failure. \\
\hline
QC-FND-005 & Omission & Medium & No glossary or definitions section; acronyms undefined within SRS. \\
\hline
QC-FND-006 & Authorship & Medium & Vague or ambiguous terms (e.g., "periodic") and non-measurable criteria. \\
\hline
QC-FND-007 & Authorship & Low & Typographical errors (duplicate "shall") and structural inconsistencies (duplicate section heading). \\
\hline
\end{tabularx}
\caption{Summary of QC findings}
\label{tab:qc-findings-summary}
\end{table}

\subsection{Detailed Findings}

\subsubsection*{QC-FND-001: Contradiction in Multi-Factor Authentication Scope}
\textbf{Related error detection:} QC-ERR-001 \\
\textbf{Type:} Contradiction \\
\textbf{Severity:} High

\textbf{Context.} Section 1.4.1 (In-Scope Deliverables) states that MFA is required for "all user accounts". Section 3.2 FR-2.6 limits MFA to "all administrative accounts".

\textbf{Result.} A developer or implementer cannot determine whether non-administrative users (e.g., suppliers, transporters, customers) must use MFA.

\textbf{Analysis.} These two statements are mutually exclusive; one must be incorrect or the scope must be clarified. This is a latent contradiction that will lead to inconsistent implementation unless resolved.

\subsubsection*{QC-FND-002: Contradiction in Blockchain Transaction Confirmation Time}
\textbf{Related error detection:} QC-ERR-008 \\
\textbf{Type:} Contradiction \\
\textbf{Severity:} High

\textbf{Context.} Section 1.4.1 (Performance Requirements) specifies blockchain transaction confirmation as "<30 seconds for standard transactions". Section 5.1 NFR-2 requires confirmation "within 10 seconds for standard operations".

\textbf{Result.} Two different performance targets exist for the same operation, creating ambiguity for design and testing.

\textbf{Analysis.} The SRS provides conflicting non-functional requirements. The customer must decide which target is correct.

\subsubsection*{QC-FND-003: Contradiction in Recovery Objectives}
\textbf{Related error detection:} QC-ERR-011 \\
\textbf{Type:} Contradiction \\
\textbf{Severity:} High

\textbf{Context.} NFR-19 (Section 5.4) sets RPO $\leq$ 24 hours and RTO $\leq$ 4 hours. CM-5 (Section 6.2.5) sets RTO $\leq$ 2 hours and RPO $\leq$ 30 minutes.

\textbf{Result.} These values are significantly different and cannot both be true.

\textbf{Analysis.} The SRS contains contradictory disaster recovery targets, which will confuse architects and operators.

\subsubsection*{QC-FND-004: Omission of Critical Exception and Edge Cases}
\textbf{Related error detection:} QC-ERR-003, QC-ERR-005, QC-ERR-007, QC-ERR-010 \\
\textbf{Type:} Omission \\
\textbf{Severity:} High

\textbf{Context.} The SRS specifies normal operational behaviour but omits several important exception scenarios:
\begin{itemize}
\item No password reset/recovery process despite account lockout.
\item No handling of blockchain unavailability or transaction failure.
\item No behaviour defined for invalid or tampered QR codes during product verification.
\item No detailed requirements for data migration validation or failure handling.
\end{itemize}

\textbf{Result.} Downstream design and implementation may fill these gaps with assumptions, leading to inconsistent or incorrect system behaviour in edge cases.

\textbf{Analysis.} These omissions represent significant completeness risks, particularly in a regulatory and safety-critical domain.

\subsubsection*{QC-FND-005: Missing Glossary}
\textbf{Related error detection:} QC-ERR-012 \\
\textbf{Type:} Omission \\
\textbf{Severity:} Medium

\textbf{Context.} The SRS contains numerous acronyms (HSA, MOH, PDPC, IMDA, GMP, RPO, RTO, etc.) but provides no glossary or definitions section.

\textbf{Result.} Readers unfamiliar with the domain may misinterpret requirements.

\textbf{Analysis.} Including a glossary is a standard practice to ensure clarity and shared understanding.

\subsubsection*{QC-FND-006: Ambiguous and Non-Measurable Requirements}
\textbf{Related error detection:} QC-ERR-004, QC-ERR-009 \\
\textbf{Type:} Authorship \\
\textbf{Severity:} Medium

\textbf{Context.} Several requirements use vague or undefined terms:
\begin{itemize}
\item FR-3.9: "periodic GPS updates" -- frequency not specified.
\item QF-3 Measurable Criteria: "100\% enforcement of role-based access control" -- lacks testable definition.
\end{itemize}

\textbf{Result.} These requirements cannot be verified objectively and may be interpreted differently by different stakeholders.

\textbf{Analysis.} The intended meaning is clear to the authors, but the written statements lack precision needed for implementation and testing.

\subsubsection*{QC-FND-007: Editorial and Structural Inconsistencies}
\textbf{Related error detection:} QC-ERR-002, QC-ERR-006 \\
\textbf{Type:} Authorship \\
\textbf{Severity:} Low

\textbf{Context.} Minor errors reduce document professionalism:
\begin{itemize}
\item Duplicate "shall" in NFR-4 and NFR-7.
\item Duplicate section heading 1.4.3 appearing twice with different content.
\end{itemize}

\textbf{Result.} While these do not affect requirement meaning, they detract from readability and suggest insufficient editorial review.

\textbf{Analysis.} Such errors are typical of early-stage documents and can be easily corrected.

% =====================
% 7 CONCLUSION
% =====================
\section{Conclusion}

The QC audit concludes that the Software Requirements Specification is substantively mature and reflects strong work carried out during the earlier Requirements Engineering stages. However, the document still contains defects typical of the validating stage, particularly in the areas of contradictions (MFA scope, performance targets, recovery objectives), omissions (exception handling, recovery processes, glossary), and authorship clarity (vague terms, editorial errors).

The QC team does not directly modify the SRS. Instead, the findings documented in this report are presented for customer consideration. The customer will determine what actions, if any, should be taken to address the identified defects.

Overall, the SRS is assessed as \textbf{conditionally acceptable}, pending customer review and resolution of the findings.

% =====================
% 8 SIGN-OFF
% =====================
\section{Sign-off}

\subsection{QC Audit Team Authentication}

\begin{table}[H]
\centering
\small
\renewcommand{\arraystretch}{1.6}
\begin{tabular}{|m{0.32\textwidth}|m{0.23\textwidth}|m{0.3\textwidth}|}
\hline
\rowcolor{tableheader}
\textbf{Name} & \textbf{Role} & \textbf{Authentication} \\
\hline
Muhammad Ridwan Putra Jasni & QC Audit Lead & Submitted by named team member \\
\hline
Lin Yuhao & QC Auditor & Submitted by named team member \\
\hline
Muhammad Mikhail Bin Mazlan & QC Auditor & Submitted by named team member \\
\hline
Kannan s/o Rajamohan & QC Auditor & Submitted by named team member \\
\hline
\end{tabular}
\caption{QC audit team sign-off}
\label{tab:qc-signoff}
\end{table}

% =====================
% 9 ADDENDA
% =====================
\section{Addenda}

\subsection{Repository Information}

\begin{center}
\fbox{\begin{minipage}{0.9\textwidth}
\centering
\textbf{PROJECT REPOSITORY}\\[2pt]
\href{https://github.com/Spooky312/The-Abyss-Navigation-Suite-INF2113}{\uline{\nolinkurl{https://github.com/Spooky312/The-Abyss-Navigation-Suite-INF2113}}}\\[4pt]
\textit{QC report, SRS, models, and traceability artefacts maintained under version control}
\end{minipage}}
\end{center}

\subsection{Glossary}

\begin{table}[H]
\centering
\small
\renewcommand{\arraystretch}{1.4}
\begin{tabularx}{\textwidth}{|p{0.28\textwidth}|X|}
\hline
\rowcolor{tableheader}
\textbf{Term} & \textbf{Definition} \\
\hline
QC & Quality Control inspection of a work product, in this case the SRS. \\
\hline
SRS & Software Requirements Specification used as the requirements baseline for downstream development. \\
\hline
Authorship defect & Defect caused by unclear or imprecise expression of intended requirement meaning. \\
\hline
Omission defect & Defect in which required content, scenario, or supporting information is missing. \\
\hline
Contradiction defect & Defect in which different parts of the SRS imply inconsistent or conflicting behaviour. \\
\hline
Traceability coherence & The degree to which requirements, models, sources, and references link together clearly and consistently. \\
\hline
MFA & Multi-Factor Authentication. \\
\hline
RPO & Recovery Point Objective -- maximum acceptable age of files restored from backup. \\
\hline
RTO & Recovery Time Objective -- maximum acceptable downtime. \\
\hline
\end{tabularx}
\caption{QC glossary addendum}
\label{tab:qc-glossary}
\end{table}

% =====================
% 10 AI REVIEW
% =====================
\section{AI Review}

The following prompts were used with an AI assistant during the QC audit process to assist in identifying potential defects. The responses provided by the AI were considered alongside manual inspection.

\subsection*{Prompt 1}
\textit{"Review the attached SRS for contradictions in security requirements, particularly around multi-factor authentication (MFA)."}

\textbf{AI Response:}
\begin{quote}
In Section 1.4.1 (In-Scope Deliverables), under "Authentication \& Security", it states: "Multi-factor authentication (MFA) for all user accounts with SMS/email OTP support." However, Section 3.2 FR-2.6 says: "The system shall enforce multi-factor authentication (MFA) for all administrative accounts." This is a direct contradiction: the first requires MFA for all users, the second only for administrators. This inconsistency needs resolution.
\end{quote}

\subsection*{Prompt 2}
\textit{"Identify any missing exception-handling requirements in the blockchain transaction recording functional requirements."}

\textbf{AI Response:}
\begin{quote}
FR-1.1 to FR-1.4 describe recording transactions, validating hash integrity, executing smart contracts, and viewing records. There is no requirement addressing failures such as blockchain network unavailability, transaction timeout, consensus failure, or retry logic. These omissions could lead to undefined system behavior in degraded modes.
\end{quote}

\subsection*{Prompt 3}
\textit{"Point out any vague or undefined terms in the functional requirements."}

\textbf{AI Response:}
\begin{quote}
FR-3.9 states: "The system shall allow Transporters to submit periodic GPS updates during shipment transit." The term "periodic" is not defined; it could mean every minute, every hour, or upon specific events. This ambiguity should be resolved by specifying a frequency or trigger condition.
\end{quote}

\subsection*{Prompt 4}
\textit{"Check for inconsistencies in performance targets across the document."}

\textbf{AI Response:}
\begin{quote}
Section 1.4.1 (Performance Requirements) states blockchain transaction confirmation should be "<30 seconds for standard transactions". Section 5.1 NFR-2 requires confirmation "within 10 seconds for standard operations". These are conflicting values and must be aligned.
\end{quote}

\subsection*{Prompt 5}
\textit{"Identify any typographical or formatting errors."}

\textbf{AI Response:}
\begin{quote}
In Section 5.1, NFR-4 and NFR-7 both contain a duplicated word: "shall shall complete". Also, Section 1.4.3 appears twice (pages 9 and 10) with different content, indicating a numbering error.
\end{quote}

% =====================
% REFERENCES
% =====================
\clearpage
\pagenumbering{arabic}
\setcounter{page}{1}
\renewcommand{\thepage}{R-\arabic{page}}
\section*{References}
\addcontentsline{toc}{section}{References}
\begin{itemize}[leftmargin=*,nosep]
\item Medicare Plus, \textit{Project Specification: Smart Contract Enabled Pharmaceutical Supply Chain System}, ICT2113, 2026.
\item Brigman Deep Sea Technology, \textit{Software Requirements Specification (SRS): Smart Contract Enabled Pharmaceutical Supply Chain System}, 2026.
\item Kevin Anthony Jones, \textit{Labtorial Guide for Class Project: Reqrs now \& near future}, Dec 2025.
\item Kevin Anthony Jones, \textit{Rubric for Assessment of Validating Report}, Mar 2026.
\item Karl E. Wiegers, \textit{Software Requirements Specification Template}, 1999.
\end{itemize}

\end{document}
